\section{What a zero determinant can reveal}

A determinant equal to zero is often easier to understand than a general determinant value. In small dimensions we can connect the equation \(\det A=0\) directly with relations among rows or columns.

\subsection{The \(2\times2\) case}

Let
\[
A=\begin{pmatrix}a&b\\c&d\end{pmatrix}.
\]
Then
\[
\det A=0
\quad\Longleftrightarrow\quad
ad=bc.
\]
If the first row is non-zero, this equation implies that the second row is proportional to the first. To see this, suppose first that \(a\ne0\). Then
\[
d=\frac{bc}{a},
\]
so
\[
(c,d)=\frac{c}{a}(a,b).
\]
If \(a=0\), then \(b\ne0\), and the equation \(ad=bc\) gives \(c=0\). Hence
\[
(c,d)=\frac{d}{b}(0,b).
\]
If the first row is the zero row, proportionality is immediate.

Thus, for \(2\times2\) matrices, determinant zero is equivalent to the rows being proportional. The same is true for the columns.

\subsection{A first observation from matrix--vector multiplication}

The previous chapter introduced multiplication of a matrix by a column vector. Consider a matrix whose second row is \(\lambda\) times its first row:
\[
A=
\begin{pmatrix}
a&b\\
\lambda a&\lambda b
\end{pmatrix}.
\]
For an arbitrary column vector
\[
x=\begin{pmatrix}u\\v\end{pmatrix},
\]
we obtain
\[
Ax=
\begin{pmatrix}
au+bv\\
\lambda au+\lambda bv
\end{pmatrix}
=
\begin{pmatrix}
t\\
\lambda t
\end{pmatrix},
\qquad t=au+bv.
\]
No matter which values of \(u\) and \(v\) are chosen, the two entries of \(Ax\) satisfy the fixed relation
\[
(Ax)_2=\lambda(Ax)_1.
\]

\begin{examplebox}
Let
\[
A=\begin{pmatrix}1&2\\2&4\end{pmatrix}.
\]
Then \(\det A=0\), and for every \(x=(u,v)^{\mathsf T}\),
\[
Ax=
\begin{pmatrix}u+2v\\2u+4v\end{pmatrix}
=
\begin{pmatrix}t\\2t\end{pmatrix}.
\]
The resulting columns are restricted to a single pattern: the second entry is always twice the first.
\end{examplebox}

This does not yet require a general theory of linear mappings. It is simply an algebraic observation obtained from matrix multiplication. It gives a first indication that determinant zero records a restriction in the possible outputs.

A similar phenomenon appears for three rows. If
\[
r_3=r_1+r_2,
\]
then for every column vector \(x\), the third entry of \(Ax\) is the sum of the first two entries. The output again satisfies a fixed relation.

\subsection{Parameters}

When a matrix contains a parameter, its determinant is an expression in that parameter. The equation
\[
\det A(t)=0
\]
identifies the parameter values for which a row or column relation appears.

\begin{examplebox}
Let
\[
A(t)=\begin{pmatrix}t&1\\4&t\end{pmatrix}.
\]
Then
\[
\det A(t)=t^2-4=(t-2)(t+2).
\]
Therefore the determinant is zero for
\[
t=2\qquad\text{or}\qquad t=-2.
\]
For \(t=2\),
\[
A(2)=\begin{pmatrix}2&1\\4&2\end{pmatrix},
\]
and the second row is twice the first. For \(t=-2\),
\[
A(-2)=\begin{pmatrix}-2&1\\4&-2\end{pmatrix},
\]
and the second row is \(-2\) times the first.
\end{examplebox}

Here the roots of the determinant are not merely numbers obtained from an equation. They identify exactly when proportionality occurs.

\begin{examplebox}
Consider
\[
B(s)=
\begin{pmatrix}
1&0&1\\
0&1&1\\
1&s&2
\end{pmatrix}.
\]
Using the \(3\times3\) formula,
\[
\det B(s)=1-s.
\]
Thus \(\det B(s)=0\) when \(s=1\). At that value,
\[
(1,1,2)=(1,0,1)+(0,1,1),
\]
so the third row is the sum of the first two rows.
\end{examplebox}

\begin{summarybox}
When a determinant is zero, first look for a relation among rows or columns. With parameters, solve \(\det A=0\) and then inspect what relation appears at the exceptional values.
\end{summarybox}
